\section{Conditional flow matching as a generic representation of posteriors}
\label{sec:methods}

This section introduces the instrument used throughout the paper. We do
\emph{not} propose a new assimilation scheme. We use the conditional-flow
representation because it is generic enough to contain every method compared
here as a special case, which turns ``method A is better than method B'' into
the sharper question of \emph{which components of the posterior each method is
able to represent at all}.

\subsection{The conditional flow}

Let $\tau \in [0,1]$ be an abstract transport variable, distinct from physical
time $t$. Define a linear probability path from a Gaussian base sample
$\xv_0 \sim \mathcal{N}(0,\sigma_0^2 \mathbb{I})$ to a target
$\xv_1 \sim p(\xv_1 \mid \Cv)$:
\begin{equation}
  \label{eq:interpolant}
  \xv_\tau = \alpha_\tau \xv_0 + \beta_\tau \xv_1 ,
  \qquad \alpha_0 = 1,\ \beta_0 = 0,\ \alpha_1 = 0,\ \beta_1 = 1 ,
\end{equation}
with the optimal-transport schedule $\alpha_\tau = 1-\tau$,
$\beta_\tau = \tau$. Differentiating~\eqref{eq:interpolant} gives an ODE whose
flow generates the target posterior:
\begin{equation}
  \label{eq:ode}
  \frac{\mathrm{d}\xv_\tau}{\mathrm{d}\tau}
  = \frac{\dot\alpha_\tau}{\alpha_\tau}\,\xv_\tau
  + \left(\dot\beta_\tau - \beta_\tau\frac{\dot\alpha_\tau}{\alpha_\tau}\right)
    \Psiop\!\left(\xv_\tau, \Cv, \tau\right),
  \qquad
  \Psiop\!\left(\xv_\tau, \Cv, \tau\right) \;\triangleq\; \mathbb{E}\!\left[\xv_1 \mid \xv_\tau, \Cv\right].
\end{equation}

Everything method-specific is carried by the single operator $\Psiop$: the
posterior-mean estimate of the target given a partially noised state and the
context. \textbf{$\Psiop$ is the object this paper measures.}

\subsection{The partition}
\label{sec:partition}

$\Psiop$ admits an exact additive decomposition,
\begin{equation}
  \label{eq:partition}
  \Psiop\!\left(\xv_\tau, \Cv, \tau\right)
  = \underbrace{\Psimean(\Cv)}_{\text{conditional mean}}
  + \underbrace{\PsiG(\xv_\tau, \Cv, \tau)}_{\text{affine gain}}
  + \underbrace{\PsiNG(\xv_\tau, \Cv, \tau)}_{\text{non-Gaussian residual}},
\end{equation}
with $\Psimean(\Cv) = \mathbb{E}[\xv_1 \mid \Cv]$ independent of $\xv_\tau$,
$\PsiG$ the $L^2$ projection of $\Psiop - \Psimean$ onto functions affine in
$\xv_\tau$, and $\PsiNG$ the remainder. Defining $\PsiG$ as an $L^2$ projection
makes the three terms orthogonal, so the decomposition is unique and the
attribution in Section~\ref{sec:results} is well posed.

\paragraph{The non-Gaussian branch vanishes at both ends.}
At $\tau=0$, $\xv_\tau \perp \xv_1$, so $\Psiop = \mathbb{E}[\xv_1\mid\Cv] = \Psimean$
and the affine gain is zero; hence $\PsiNG(0)=0$. At $\tau=1$, $\xv_\tau = \xv_1$
and $\Psimean + \PsiG = \xv_1$ exactly; hence $\PsiNG(1)=0$. The non-Gaussian
term is therefore a \emph{bump} on $(0,1)$, pinned to zero at both boundaries.

\subsection{What each method class can represent}
\label{sec:zeroing}

The value of~\eqref{eq:partition} as a diagnostic is that the classical schemes
are not merely different points in a hypothesis space --- they \emph{identically
zero} named terms.

\begin{table}[t]
\centering
\small
\begin{tabular}{@{}lccc@{}}
\toprule
\textbf{method class} & $\Psimean$ & $\PsiG$ & $\PsiNG$ \\
\midrule
Optimal interpolation, 4D-Var & $\hat{\xv}(\Cv)$, constant in $\xv_\tau$ & $\mathbf{0}$ & $\mathbf{0}$ \\
DirectUNet (deterministic regression) & $\hat{\xv}(\Cv)$, constant in $\xv_\tau$ & $\mathbf{0}$ & $\mathbf{0}$ \\
EnKF / ETKF / ensemble smoothers & ensemble mean & \checkmark\ (the Kalman gain \emph{is} $K_\tau$) & $\mathbf{0}$ \emph{by construction} \\
CFM, score-based assimilation (SDA) & \checkmark & \checkmark & \checkmark \\
DirectUNet\,$+$\,SDA (blending) & imported, deterministic & \checkmark & \checkmark \\
\bottomrule
\end{tabular}
\caption{What each class of scheme can represent. The zeros are structural, not
symptoms of tuning: a Gaussian update \emph{is} the ensemble-Kalman method, and a
deterministic estimator ignores its state argument by definition.}
\label{tab:zeroing}
\end{table}

Two consequences frame the rest of the paper.

\paragraph{A deterministic estimator is a degenerate flow.}
Read any deterministic estimator $\hat{\xv}(\Cv)$ as an operator that ignores its
state argument, $\Psiop_{\det}(\xv_\tau,\Cv) = \hat{\xv}(\Cv)$. Substituting
into~\eqref{eq:ode} and integrating gives exactly
$\xv_\tau = (1-\tau)\,\xv_0 + \tau\,\hat{\xv}$. 4D-Var, optimal interpolation and
a supervised regression network are therefore \emph{the same member} of this
family, differing only in how $\hat{\xv}$ is computed. This is why the paper can
place them on one axis without apology.

\paragraph{The ensemble-Kalman class has a structural ceiling.}
$\PsiNG \equiv 0$ is not a deficiency of a particular implementation or
inflation setting: it is what it means to be an ensemble-Kalman method. Any
posterior feature living in $\PsiNG$ is therefore unreachable for that class at
\emph{any} tuning --- which is the precise form of symptom S-b, and what
Section~\ref{sec:nongaussian} measures.

\paragraph{Blending composes the slots.}
Warm-starting a generative sampler from a deterministic estimate --- the
DirectUNet\,$+$\,SDA schemes evaluated here --- is not an unprincipled trick. In
this representation it is the statement that $\Psimean$ may be supplied by one
scheme while $\PsiG$ and $\PsiNG$ are supplied by another. Section~\ref{sec:results}
measures what that composition buys and what it costs.
